\begin{claim}{AG.35}{Proved: geometric origin alone does not bound unit-root returns}
Consider the fixed smooth elliptic family over the multiplicative group
\[
 \mathcal E\longrightarrow\mathbf G_{m,\mathbf F_4},
 \qquad E_a:\ y^2+xy=x^3+a.                         \tag{AG.35.1}
\]
For every positive odd integer $m$, there is a closed parameter $a$
of exact degree $m$ over $\mathbf F_4$ such that
\[
 \#E_a(\mathbf F_{4^m})=4^m,
 \qquad E_a(\mathbf F_{4^m})\simeq\mathbf Z/2^{2m}\mathbf Z.
                                                        \tag{AG.35.2}
\]
If $\beta_a\in\mathbf Z_2^\times$ is its unit-root Frobenius
eigenvalue, then
\[
 \beta_a^2-\beta_a+4^m=0,
 \qquad v_2(\beta_a-1)=2m.                         \tag{AG.35.3}
\]
The parameter is a point of odd order in
$\mathbf G_m(\overline{\mathbf F}_2)$, and its Teichmuller lift
is prime-to-two torsion. Thus a single fixed geometric family,
ordinary fibers, integral cohomology of Weil weight one, and
prime-to-two Frobenius periods do not give a uniform bound for
the return unit minus one, even at prime-to-two torsion parameters.
The assertion does not put these parameters in the powers of one
fixed odd prime, and does not identify this family with the Abel
tower of AR.20. Those restrictions remain essential unresolved
inputs to that route.
\end{claim}
\begin{proof}
The discriminant and $j$-invariant of (AG.35.1) are $a$ and $1/a$,
respectively. Hence the displayed projective Weierstrass model is
a smooth elliptic scheme over $\mathbf G_m$. Each fiber is ordinary:
the geometric point $(0,\sqrt a)$ is a nonzero point of order two,
whereas an elliptic curve has at most one independent geometric
two-torsion point in characteristic two.

Fix odd $m$ and put $q=4^m$. Waterhouse's classification
\citep[Theorem~4.1, case~(1), p.~536]{AG_Waterhouse1969}
asserts that every integer $t$ with $|t|\le2\sqrt q$ and
$(t,2)=1$ is the Frobenius trace of an ordinary elliptic curve
over $\mathbf F_q$. Apply it to $t=1$. The resulting elliptic
curve $E/\mathbf F_q$ has $q$ rational points.
Since $4\mid q$, Proposition~5.1 of
\citep[p.~311]{AG_AuerTop2002} identifies $E$ over $\mathbf F_q$
with a member $E_a$ of (AG.35.1).

For completeness, the normal-form step can be seen directly.
An ordinary elliptic curve over the perfect field $\mathbf F_q$
has an equation
\[
               y^2+xy=x^3+b x^2+a,\qquad a\ne0.
\]
Its unique nonzero geometric two-torsion point is
$T=(0,\sqrt a)$. As $E(\mathbf F_q)$ is a two-group and
the geometric two-primary torsion has rank one, it is cyclic;
since its order is $q\ge4$, the point $T$ has a rational half.
The duplication formula for a point $R$ with nonzero
$x$-coordinate is
\[
                 x(2R)=x(R)^2+\frac a{x(R)^2}.
\]
Thus a half of $T$ must have $x$-coordinate $r=a^{1/4}$.
Writing its $y$-coordinate as $rz$, the equation becomes
\[
                       z^2+z=b+r+r^2.
\]
It is solvable over $\mathbf F_q$ precisely when
$\operatorname{Tr}_{\mathbf F_q/\mathbf F_2}(b)=0$.
Consequently $b=s^2+s$ for some $s\in\mathbf F_q$, and
the substitution $y\mapsto y+sx$ removes the $x^2$ term.
This proves the required isomorphism to $E_a$ and also confirms
the rational four-torsion step in the cited proposition.

The characteristic polynomial of $E_a$ is $X^2-X+q$.
Modulo two it has distinct roots zero and one. Its unit root
$\beta_a$ therefore belongs to $\mathbf Z_2^\times$, and the
other root is $q/\beta_a$, of valuation $2m$. Evaluating at one
gives
\[
              (1-\beta_a)(1-q/\beta_a)=q.
\]
The second factor is a two-adic unit, proving (AG.35.3).
The geometric two-primary torsion of an ordinary elliptic curve
is $\mathbf Q_2/\mathbf Z_2$, so its finite rational two-primary
subgroup is cyclic. Its full rational order here is $q$, which
proves the group assertion in (AG.35.2).

It remains to prove that the parameter has exact degree $m$;
the existence of a curve over $\mathbf F_{4^m}$ alone would not
establish the assertion about periods. Let
$d=[\mathbf F_4(a):\mathbf F_4]$, so $d\mid m$. The displayed
model $E_a$ is defined over $\mathbf F_{4^d}$. Write
$\gamma\in\mathbf Z_2^\times$ for its unit root over that field.
The unit root over $\mathbf F_{4^m}$ is
$\gamma^{m/d}$. Because $m/d$ is odd,
\[
 v_2(1-\gamma^{m/d})=v_2(1-\gamma);
\]
the quotient $(1-\gamma^{m/d})/(1-\gamma)$ is a sum of an
odd number of odd two-adic integers. The nonunit Frobenius root
does not contribute to the two-primary part of the rational
point count, so
\[
             v_2\bigl(\#E_a(\mathbf F_{4^d})\bigr)=2m.
\]
If $d<m$, Hasse's bound contradicts this divisibility:
\[
 4^m\le\#E_a(\mathbf F_{4^d})
       \le(2^d+1)^2<4^m.
\]
Hence $d=m$. Finally $a\in\mathbf F_{4^m}^\times$ has order
dividing the odd number $4^m-1$, and its fourth-power Frobenius
orbit has length exactly $m$. Teichmuller lifting preserves that
order and that orbit length.

The small exact certificate
\path{agent_geometric/code/geometric_unitroot_growth.py}
constructs examples for $m=1,3,5$, records the finite fields,
the parameter, its order and degree, and the point count, and
checks the unit-root congruence. This finite check illustrates
the theorem; its assertion for every odd $m$ follows from the
classification and the degree argument above.
\end{proof}
